\chapter*{Referencial Teórico}
\addcontentsline{toc}{chapter}{Referencial Teórico}

Este trabalho fundamenta-se nos conceitos de sistemas ciberfísicos, robótica autônoma e visão computacional embarcada, no contexto da Quarta Revolução Industrial, caracterizada pela integração entre os domínios físico e digital \citeonline{Schwab2016}. A visão computacional é compreendida como o processo de inferência de propriedades do mundo físico a partir de imagens, permitindo a percepção e interpretação do ambiente por sistemas artificiais \citeonline{Szeliski2021}. A autonomia robótica emerge da integração entre percepção, tomada de decisão baseada em inteligência artificial e controle cinemático, operando de forma acoplada e em tempo real. Para viabilizar essa integração em plataformas de baixo custo, adota-se o paradigma da Computação de Borda (\textit{Edge Computing}), que prioriza o processamento local dos dados, reduzindo latência e dependência de infraestrutura externa \citeonline{BritoGuimaraes2023}. Esses fundamentos orientam o desenvolvimento do manipulador robótico SCARA proposto, concebido como uma plataforma autônoma capaz de perceber, decidir e agir no ambiente físico.